\section{Background}

\subsection{CXL Fabric-Attached Memory}

CXL fabric-attached memory (FAM) enables multiple hosts to transparently share a common memory pool, reducing data movement overhead for AI training, inference serving, and data analytics~\cite{tangYggdrasilReducingNetwork2024,zhongMyCXLPool2025,liPondCXLBasedMemory2023,goukDirectAccessHighPerformance}.
Sharing memory demands synchronization: parallel communication primitives (AllReduce, AllGather) require scalable locks~\cite{dangAdvancedThreadSynchronization2017}, and non-scalable locks can cause dramatic performance collapse~\cite{boyd-wickizerNonscalableLocksAre}.
The choice of synchronization mechanism---and the coherence model it requires---is a first-order architectural decision.

\subsection{Coherence Models in CXL}

CXL supports cache coherence from its earliest specification.
CXL.cache provides host-device coherence for Type~1 devices; Type~2 devices add CXL.mem with host-managed device memory and device-managed coherence (HDM-DB), introducing the back-invalidation snoop (BISnp) protocol~\cite{sharmaIntroductionComputeExpress}.
Type~3 devices with Multi-Logical Devices (MLDs) enable shared FAM by exposing memory to multiple hosts over port-based routing.

For multi-host shared FAM, the CXL 3.X specification defines two coherence models:

\paragraph{Hardware-managed coherence.}
Extends BISnp across hosts, maintaining an inclusive snoop filter that tracks cache-line ownership per host.
When a host requests a cache line already held by another host, the device issues a back-invalidation (BI) message to the current owner, forcing it to writeback and relinquish the line before granting ownership to the requester.
This protocol enables RMW atomic operations: the device grants exclusive cache-line ownership to a requesting host, which then executes the atomic locally within its cache hierarchy.
The cost is threefold: (1)~device-side snoop filter storage proportional to the number of tracked cache lines and hosts, (2)~cross-host invalidation round-trips that add latency to every ownership transition, and (3)~protocol complexity that scales with host count~\cite{goukDirectAccessHighPerformance}.
For synchronization, this cost manifests directly: every contested lock acquisition requires at least one ownership transition, incurring the full BI round-trip latency.

\paragraph{Software-managed coherence.}
Avoids hardware coherence entirely; the software layer manages shared state through explicit loads, stores, and fences.
This eliminates coherence infrastructure but restricts synchronization to software locks.
A key concern is that CXL 3.X's support for atomic loads/stores under software-managed coherence is ambiguous~\cite{lamportMutualExclusionProblem,lamportMutualExclusionProblema}: simultaneous multiport read/write operations can cause flickered reads or scrambled writes, a risk that grows as devices support up to 16 ports.
Support for atomic loads/stores does not fundamentally require full cache coherence, as conflict resolution schemes exist~\cite{wangIntelligentMultiPortMemory2008}, but whether CXL 3.X devices will implement such schemes remains unknown.

\subsection{Memory Access Modes for CXL}
\label{sec:access-modes}

Beyond the coherence model, the \emph{memory access mode} determines how the host CPU interacts with CXL-attached memory.
These access modes are configured via page table attributes (specifically the PAT, PCD, and PWT bits on x86-64) and are orthogonal to the device-side coherence model.
We identify four modes relevant to synchronization:

\paragraph{Cached with hardware coherence (Cached-HC).}
The default mode when hardware coherence is available.
Loads and stores go through the host cache hierarchy; the BISnp protocol ensures cross-host consistency.
RMW atomics are supported.
This mode benefits from cache locality---repeated reads to the same lock variable hit L1/L2 when no other host has written---but incurs coherence traffic proportional to the number of sharing hosts.
The performance of synchronization under this mode is bounded by the coherence protocol's ownership-transition latency.

\paragraph{Cached with software coherence (Cached-SC).}
Used when the device provides only software-managed coherence.
Host caches are active, but consistency is the software's responsibility via explicit fences and flushes.
No RMW atomics are available across hosts; only software locks are viable.
The challenge is ensuring that stores to lock variables become visible to other hosts promptly; this requires \texttt{clflush} or \texttt{clwb} instructions after stores, followed by \texttt{sfence} to order the flush.
The cost of these explicit consistency operations partially offsets the benefit of cached access.

\paragraph{Uncached (UC).}
The host page table marks CXL memory as uncacheable (PAT entry set to UC).
Every load and store traverses the CXL link directly, bypassing the host cache hierarchy entirely.
This eliminates all coherence traffic---no snoop filter entries, no invalidation messages---at the cost of losing cache-line spatial and temporal locality.
Each memory access pays the full CXL round-trip latency.
On x86-64, UC accesses are strongly ordered: no additional fences are needed, as the processor guarantees that UC loads and stores are not reordered.
For synchronization, this means that polling a lock variable incurs a CXL round-trip on every iteration, but there is zero coherence overhead, no risk of stale cached copies, and no need for explicit flush instructions.

\paragraph{Non-temporal (NT).}
Uses streaming instructions (\texttt{movntdq} for 128-bit stores, \texttt{movntdqa} for 128-bit loads on SSE4.1+ platforms) that bypass the cache hierarchy.
NT stores exploit write-combining (WC) buffers: the processor coalesces multiple stores to the same WC buffer entry (typically 64 bytes, matching a cache line) into a single bus transaction.
This reduces CXL link utilization for write-heavy paths by up to $8\times$ when a full cache line of stores can be combined.
NT loads avoid cache pollution but do not benefit from write-combining.
For lock implementations, NT access is most beneficial on the unlock path where a single store must be globally visible, as the write-combining buffer can merge the store with any adjacent metadata updates.
A key constraint is that NT stores require an \texttt{sfence} to ensure global visibility, adding a serialization point after each store.

\paragraph{Why access modes matter for synchronization.}
Prior work on CXL synchronization~\cite{suetterleinSynchronizationCXLBased2024} implicitly assumes cached access, leaving uncached and non-temporal modes entirely unexplored.
Yet these modes represent the simplest deployment options: UC access requires no coherence support at all, and NT access further optimizes write traffic.
If mutex performance under UC or NT access is ``good enough'' for a given workload, the deployment can avoid the cost and complexity of hardware coherence infrastructure entirely.
This is a significant practical consideration: hardware coherence requires snoop filter SRAM on the device (scaling linearly with tracked cache lines), additional protocol state machines, and validation against the CXL specification's coherence requirements.
Quantifying ``good enough'' is a central goal of this paper.


\subsection{Mutual Exclusion}

Mutual exclusion synchronizes concurrent threads by ensuring only one thread enters a critical section at a time.
Mutex performance depends on two factors specific to CXL FAM: the availability of RMW atomics (which require hardware coherence) and the notification mechanism (limited to polling, as no centralized kernel monitor exists across hosts).

\subsubsection{Polling as Notification Mechanism}

In CXL FAM, no centralized kernel monitor can identify lock state across hosts, so all locks must rely on polling.
The cost of polling depends on the access mode:
under Cached-HC, repeated polls hit the local cache (zero CXL traffic);
under Cached-SC, polls hit the cache but require periodic flushes to detect remote writes;
under UC, every poll generates a CXL round-trip;
under NT, polls use streaming loads that bypass the cache.

\subsubsection{Software Locks}

Software locks rely on assignment and control structures rather than RMW atomics.
Following Lamport's classification~\cite{lamportMutualExclusionProblem,lamportMutualExclusionProblema}, they are either \emph{RW-safe} (correct without atomic loads/stores) or \emph{RW-unsafe} (requiring atomic loads/stores).
This distinction matters for CXL software-managed coherence where atomic load/store support is ambiguous.

Previous work~\cite{suetterleinSynchronizationCXLBased2024} addressed this with a Tournament Peterson lock that constrains shared variables to binary values, mitigating flicker/scramble risks.
We evaluate both RW-safe and RW-unsafe locks (Table~\ref{tab:Locks}) to characterize this trade-off.

Software locks must also contend with instruction reordering by compilers and hardware.
We instrument each implementation with the minimum fences required for correctness.
A notable constraint of software locks is their storage overhead: each lock requires $O(N)$ to $O(N \log N)$ bytes where $N$ is the maximum thread count, as each thread needs dedicated state variables (e.g., choosing/number arrays in Bakery, flags in Peterson).
This contrasts with hardware locks that require $O(1)$ storage for the lock word itself.

\subsubsection{Hardware Locks}

Hardware locks leverage RMW atomic instructions (test-and-set, fetch-and-add, compare-and-swap) to establish mutual exclusion.
These atomics require hardware cache coherence to grant exclusive cache-line ownership.

The performance of hardware locks in CXL environments is fundamentally constrained by two factors: (1)~the latency of the coherence protocol to acquire exclusive ownership, and (2)~cache-line contention when multiple hosts poll a shared variable.
We hypothesize that these costs, not the relative performance of different atomic instructions, dominate lock performance in FAM environments.

Not all atomic instructions are equivalent: Herlihy's consensus hierarchy~\cite{herlihyWaitfreeSynchronization1991} ranks them by solving consensus for different numbers of processes, and their performance varies by hardware implementation~\cite{schweizerEvaluatingCostAtomic2020}.
We do not explore equivalent implementations of the same lock with different atomics, as we focus on the coherence protocol overhead rather than instruction-level differences.

\subsubsection{Interaction of Access Modes and Lock Types}

Not all combinations of access modes and lock types are meaningful.
Hardware locks \emph{require} cached access with hardware coherence, as RMW atomics need exclusive cache-line ownership.
Software locks can operate under any access mode, since they rely only on loads and stores:
\begin{itemize}
    \item Under \textbf{Cached-SC}, software locks benefit from cache locality but must manage consistency via fences and explicit flushes (\texttt{clflush}/\texttt{clwb} + \texttt{sfence}).
    \item Under \textbf{UC}, every poll traverses the CXL link, eliminating coherence overhead but increasing per-access latency. The strong ordering guarantees of UC access on x86-64 eliminate the need for most fences.
    \item Under \textbf{NT}, write-combining can reduce the cost of lock release (store to a shared variable), while reads still pay full CXL latency. An \texttt{sfence} is required after NT stores.
\end{itemize}
This interaction is a key dimension of our evaluation: we measure whether the elimination of coherence traffic under UC/NT access compensates for the loss of cache locality.
